% D.89 proposed candidate: notes only, not registered or integrated.
\subsection{Cusp parameters with primitive valuation at five}

\begin{statement}{D.89}{Open}{Candidate: the entire compact kernel for a primitive five-valuation cusp}
Let
\[
 a=A\kappa(5)+\sum_{q\ne5}a_q\kappa(q),\qquad
 A\in\mathbb Z_5^\times,\qquad a_q\longrightarrow0,
\]
and consider the actual cusp pair \(s=(a,0)\).
If all \(a_q\) vanish, then \(\mathcal S_s=\mathcal K_s=0\).
Otherwise put
\[
 \begin{aligned}
 r&=\min_{q\ne5}v_5(a_q),&
 \lambda&=\sum_{q\ne5}a_q\frac{\log(q^4)}{4\log6},\\
 N&=v_5(\lambda),&e&=N-r\in\mathbb Z_{\ge0}\cup\{\infty\}.
 \end{aligned}                                             \tag{D.89.1}
\]
Then the compact groups in D.63 satisfy
\[
 \mathcal S_s\simeq\mathcal K_s\simeq
 \begin{cases}
  0,&e=0,\\
  \displaystyle\prod_{j\ge0}\mathbb Z/5^e,&0<e<\infty,\\
  \displaystyle\prod_{j\ge0}\mathbb Z_5,&e=\infty.
 \end{cases}                                               \tag{D.89.2}
\]
These are isomorphisms with the actual compact topologies.
In particular the finite-depth compact kernel has no additional
free factor. The pair always has an actual full cusp section.
For every pair of integers \(r\ge0,e\ge1\), the rational
tangent
\[
                  5(1+5^{e+1})^{5^r}                     \tag{D.89.3}
\]
realizes these indices and the exact bounded exponent \(5^e\).
Its nonempty metabelian fibre has translation group
\(\prod_{\mathbb N}\mathbb Z_5\), by D.74.

The proof retains the actual \(\kappa(5)\)-graph action.
It does not extend this formula to a nonzero nonunit \(A\),
which imposes a different cyclotomic relation. It makes no
general section-existence or first higher-derived-obstruction claim.
\end{statement}

\begin{proof}[Candidate proof, pending complete audits]
Use \(\Lambda=\mathbb Z_5[[u,v]]\), \(t_1=1+u\),
\(t_2=1+v\), \(w=t_1t_2-1\), and \(P=uvw\).
Let \(I_0\) be the actual untwisted Alexander operator image,
including \(\chi\). As in D.82, write \(\kappa_5=\kappa(5)\)
for its actual continuous weight-one cocycle; this is not the
weight-five Soul\'e character. A superscript \(D\) means
Cartier dual, and a star means ordinary Pontryagin dual after
the Cartier twist is accounted for.

\paragraph{The actual groups retain their common valuation graph.}
If the away-five class is zero, D.83 applies with its zero
lattice and proves the statement. Otherwise its closed
translation lattice is \(5^r\mathbb Z_5\) on the first axis.
The exact image and integral gauge construction in D.83/D.87
gives
\[
 I_s=5^r\mathbb Z_5\rtimes I_0,
 \qquad I_5=5^N\mathbb Z_5\rtimes I_0,                  \tag{D.89.4}
\]
where \(5^\infty\mathbb Z_5=0\). In these descriptions
the common embedded complement is the actual graph
\[
             g\longmapsto(A\kappa_5(g),g).
\]
Its action on every coefficient is
\[
                \rho_A(g)=t_1^{A\kappa_5(g)}\rho_0(g).
                                                               \tag{D.89.5}
\]
The integral local gauge lies in the actual global translation
lattice. Uniform \(c_0\) convergence gives the displayed
\(\lambda\) and \(N\ge r\), including infinite support.
No independent formal action has replaced the Galois image.

If \(e=0\), the actual images in (D.89.4) agree, so D.87
gives \(\mathcal S_s=0\). We will check the compact free
term at the end. Suppose henceforth \(e>0\), finite or infinite.
Put \(q_n=(1+u)^{5^n}-1\), including \(q_0=u\), and set
\[
 B=\Lambda/(q_r),\qquad
 R=\begin{cases}\Lambda/(q_N),&N<\infty,\\
                 \Lambda,&N=\infty,\end{cases}
 \qquad J=q_rR.                                         \tag{D.89.6}
\]
The actual coinduced inclusion and D.87 give
\[
 \mathcal S_s^*\simeq
       (J^D)^{I_0}/\operatorname{im}((R^D)^{I_0}),        \tag{D.89.7}
\]
with precisely the graph action (D.89.5).

\paragraph{The operator-difference ideal is exactly \((u)\).}
Let
\[
 N_0=\ker(\chi:I_0\to\mathbb Z_5^\times),\qquad
 N_1=\{g\in N_0:\kappa_5(g)=0\}.
\]
The already proved actual statements in D.77/D.82 supply
an element \(\delta\in N_1\) with
\(\beta_\delta-1=P\varepsilon\), where
\(\varepsilon\in\Lambda^\times\), and an element
\(\tau\in N_0\) with \(\kappa_5(\tau)=1\).
For clarity, the first assertion is the integral projection
of the primitive cubic character to the weight-three part of
the compact actual image of \((\kappa_5,\kappa_3)\):
the tame weight-one part is removed, while the cubic coordinate
remains a unit. The second uses the actual saturated image of
\(\kappa(5)\) on the cyclotomic kernel proved in D.82.
Both elements are actual operator-image elements, not assigned
values of unrelated characters.

On \(N_0\), the operator (D.89.5) is multiplication by
\(t_1^{A\kappa_5(g)}\beta_g\). Every \(\beta_g-1\)
is divisible by \(P\), by D.39. The element \(\delta\)
therefore puts \(P\) in the ideal generated by the operator
differences. The element \(\tau\), modulo that ideal,
puts \(t_1^A-1\) in it, since
\[
 t_1^A\beta_\tau-1=(t_1^A-1)+t_1^A(\beta_\tau-1).
\]
Conversely, every difference belongs to
\((P,t_1^A-1)\): for each \(c\in\mathbb Z_5\),
\(t_1^{Ac}-1\) is divisible by \(t_1^A-1\) by the
integral binomial identity. As \(A\) is a unit,
\((t_1^A-1)/u\) has unit constant coefficient \(A\).
Thus the exact closed ideal is
\[
       (\rho_A(g)-1:g\in N_0)=(P,t_1^A-1)=(u).         \tag{D.89.8}
\]
The equality is integral. In fact the two specified differences
already generate it as a \(\Lambda\)-ideal, so no closure
of a merely dense image is substituted for this equality.

Consequently, for \(X=R\) and for the actual submodule
\(X=J\), invariance of the discrete dual under \(N_0\)
means exactly annihilation of \(uX\). In particular
\[
                (R^D)^{N_0}=(R/uR)^D,
                \qquad (J^D)^{N_0}=(J/uJ)^D.            \tag{D.89.9}
\]

\paragraph{The natural invariant map vanishes.}
The map in (D.89.7) is induced by the actual inclusion
\(J=q_rR\hookrightarrow R\). Since \(q_r\) is
divisible by \(u\), its image lies in \(uR\).
Every \(N_0\)-invariant functional on \(R\) annihilates
this image by (D.89.9). Therefore
\[
                 (R^D)^{I_0}\longrightarrow(J^D)^{I_0}
                           \quad\text{is zero}.        \tag{D.89.10}
\]
This is an equality of the actual maps, not a dimension or
corank comparison. Equations (D.89.7)--(D.89.10) identify
the whole compact finite-localization term with
\[
                   \mathcal S_s\simeq(J/uJ)(-1)_\Gamma,
                   \qquad\Gamma=\mathbb Z_5^\times.    \tag{D.89.11}
\]
D.87 supplies the one finite global coefficient level for
each dual class, including every local place.

\paragraph{Finite local index and the retained generator.}
When \(N<\infty\), multiplication by the actual generator
\(q_r\) identifies the scalar module underlying \(J\)
with \(\Lambda/(q_N/q_r)\). Its quotient by \(u\) is
\[
                         J/uJ\simeq(\mathbb Z/5^e)[[v]].
                                                               \tag{D.89.12}
\]
Indeed \((q_N/q_r)(0)=5^{N-r}\); the distinguished
polynomial and cancellation details are proved in D.88.
The substitution on the retained generator is
\(q_r\mapsto q_rS_\chi(t_1^{5^r})\), which contributes
\(\chi\) at \(u=0\). The actual graph multiplier in
(D.89.5) is one at \(u=0\); the untwisted Alexander
multiplier there is \(\chi S_\chi(t_2)\).
After the Cartier untwist, (D.89.12) therefore has the exact
action
\[
       f(v)\longmapsto\chi S_\chi(1+v)
                       f((1+v)^\chi-1).                \tag{D.89.13}
\]
No division by \(5^r\) or \(5^e\) occurs.

Multiplication by \(v\) identifies this module with the
augmentation ideal in
\((\mathbb Z/5^e)[[\mathbb Z_5]]\), acted on by
\(t^x\mapsto\chi t^{\chi x}\). The missing constant
quotient has tame character \(\chi\), so tame averaging
identifies their compact coinvariants. As in D.88, the
invariant discrete dual consists of continuous functions
\(f(\gamma x)=\gamma^{-1}f(x)\), one arbitrary value
in \(\mathbb Z/5^e\) on each nonzero valuation shell,
zero at the origin, and only finitely many nonzero shell
values. Its compact dual is precisely
\(\prod_{j\ge0}\mathbb Z/5^e\). This proves the middle
case of (D.89.2), including its exact exponent and topology.

\paragraph{The infinite local index is computed directly.}
If \(N=\infty\), then \(R=\Lambda\) and
\(J=q_r\Lambda\). Since \(\Lambda\) is a domain,
division by its displayed generator is an actual scalar
identification, and
\[
                            J/uJ\simeq\mathbb Z_5[[v]].
                                                               \tag{D.89.14}
\]
The same retained-generator calculation gives (D.89.13).
Now the weighted continuous shell functions take values in
\(\mathbb Q_5/\mathbb Z_5\); continuity at the origin
again gives a direct sum, and their compact dual is
\(\prod_{j\ge0}\mathbb Z_5\). Together with the exact
zero map (D.89.10), this proves the infinite case. It is not
an extrapolation from finite exponents or infinite coranks.

\paragraph{No free local factor, and actual tangents.}
The global depth and local depth at five of the nonzero
boundary class \(a\) are both zero, because its valuation
coefficient \(A\) is a unit. The same holds for \(a-b=a\).
The class \(b=0\) contributes the identical countable
product on both sides of D.77's local/global boundary map.
Thus that map is an isomorphism, and the compact exact
sequence in D.63 identifies \(\mathcal K_s=\mathcal S_s\)
topologically. This also completes the zero-index and pure
valuation cases.

For (D.89.3), take \(A=1\). Its away-five coefficient
at two is \(5^r\), since \(1+5^{e+1}\equiv2\pmod4\).
Its local unit logarithm coefficient has valuation \(r+e\),
as proved and exactly checked in D.88. The zero local unit
case is realized by adding \(\kappa(5)\) to the completed
parameter in (D.88.24). These are actual cusp parameters,
so the full cusp construction supplies sections. Each has
finite valuation support in the stated examples, permitting
the D.74 fibre computation. The infinite torsion in the
finite-depth compact kernel is a second-cohomology feature,
not torsion in the first-cohomology translation group.
\end{proof}
